\scanpage{122}
\mspage{60}
On désigne par $T^!_d$ l'image inverse de $T_d$ sur $D^!_d$, sur lequel $\mathfrak S_d$ opère donc (c'est un \deleted{\unsure{rev}.} syst. local sur $(D^!_d,\mathfrak S_d)\approxeq D_d$). Or sur $D^!_d$ on a les \inserted{$d$} morphismes
\[
(185)\qquad s_i:D^!_d\mathrel{\tikz[baseline=-.5ex,x=1em,y=1ex]{\draw (0,0)--(1.8,0)--(1.84,-.2);}}D_1\qquad(1\leq i\leq d)
\]
permutés par $\mathfrak S_d$, et d'autre part, si $D_1\mathrel{\tikz[baseline=-.5ex,x=1em,y=1ex]{\draw (0,0)--(1.8,0)--(1.84,-.2);}}D_0$ est une immersion \tikz[baseline=(reg122.base)]{\node[inner sep=0pt,outer sep=0pt] (reg122) {régulière};\path (reg122.south west)--(reg122.south east) coordinate[pos=.33] (reg122end);\draw ([xshift=-.2mm]reg122.south west)--(reg122end);} \inserted{de codim 1}, on a \deleted{le \inserted{sur $D_1$} \unsure{faisceau}.} le faisceau \uline{normal} par l'immersion
\[
(186)\qquad L_1\quad\begin{gathered}\text{faisceau inversible}\\\text{canonique sur }D_1\end{gathered}
\]
et on peut considérer sur $D^!_d$ les $d$ faisceaux inv.
\[
(187)\qquad L_{d,i}=s_i^*(L_1)\qquad1\leq i\leq d.
\]
Ils correspondent à des torseurs sous $(\mathbb G_m)_{D_d}$
\[
(188)\qquad\underline L_{d,i}^*=\begin{gathered}{\mathbb G_m}_{D_d}\text{-torseur sur }D_d\text{ associé}\\\text{au faisceau inv. }L_{d,i}\end{gathered}
\]
et on peut considérer
\[
(189)\qquad P^!_d=\prod_{1\leq i\leq d}\underline L_{d,i}^*\qquad\text{torseur sous }\mathbb G_m^d
\]
Comme $\mathfrak S_d$ opérant sur $D^!_d$ permute les \inserted{$s_i$ donc les} $\underline L_{d,i}$ entre eux, il opère sur $P_d$, \deleted{qui} \inserted{de} \deleted{devient} façon compatible avec son opération sur
\margin{121}
